The current manuscript draft is intentionally conservative. It treats the repository as the primary source of truth and avoids overstating what the present evidence can support. That conservative framing should not be read as uncertainty about everything in the paper. What the current evidence does establish is narrower but still useful: within this repository, a coupled memory--weighting--reflection loop can be serialized, rerun, and evaluated in a traceable way, and that loop produces measurable but uneven gains inside the local simulator.

First, the strongest aspect of the evidence is the repository's internal consistency. The same structured plan objects feed the generator, simulator, result serializer, and DoDAF/C2SIM exporters. This makes it possible to trace reported numbers back to stored JSON files and to inspect the logic that produced them. For a first journal draft, such traceability is more valuable than broader but weakly grounded claims.

Second, the selected evidence suggests that the interaction between memory augmentation, adaptive weighting, and reflection is useful within the current simulator, but also that the gains remain scenario-dependent. The full pipeline is best in the chosen single-scenario ablation, leading all five stages and all five headline metrics, and exceeds the pure baseline in all five scenarios of the refreshed cross-scenario suite. At the same time, the stage-level evidence shows that the gains are not perfectly uniform. As emphasized by the single-scenario stage table and stage-profile figure, \texttt{breach} remains the most persistent bottleneck (0.6201 under Full vs.\ 0.6768--0.6772 for detect and sustain) despite receiving the largest absolute improvement ($+0.0181$ over the pure baseline). The optimized HOPE SDE parameters ($\theta=0.9$, $\epsilon=0.005$) and the BottleneckDetector's dynamic boost together direct a larger share of the adaptation budget toward \texttt{disrupt} and \texttt{breach}, producing the largest Full-vs-next-best gaps in exactly those historically stubborn stages ($+0.0319$ in disrupt, $+0.0181$ in breach).

The newly introduced \texttt{BottleneckDetector} module addresses this diagnostic gap by making bottleneck identification explicit and dynamic. The detector's output confirms that breach is the active bottleneck in every iteration of the single-scenario run, with stable severity scores around 0.053--0.054. The dynamic boost vector it produces directs additional emphasis to mobility, fires, and protection---the three capabilities that feed the breach stage---in proportion to the measured severity. While the dynamic boost alone does not eliminate the breach bottleneck (the stage-score evidence shows breach still lagging), it provides a principled mechanism for the controller to concentrate its adjustment budget on the most urgent stage, replacing the original static priority ordering that always privileged disrupt and breach equally. The weight-competition analysis reveals a structural reason for breach's persistence: breach and sustain share the \texttt{protection} capability dimension with equal weight (0.25), creating a tension that a scalarized reward cannot fully resolve. This finding connects the pipeline's behavior to the broader multi-task optimization literature, where gradient conflicts between competing objectives are a known challenge \citep{yu2020pcgrad}.

The memory-conditioned Hope coupling (Section~2.9) is a further step toward tighter component interaction. By letting retrieved case metrics bias the fast-weight initialisation, the pipeline no longer treats memory and adaptive weighting as purely serial modules. The correction is intentionally lightweight (clamped to $[-0.18, 0.28]$, blended at 0.30 weight) so that it nudges rather than overrides the adapter's scenario-conditioned output. The practical effect is that when the Memento module retrieves a case that failed badly at breach, the controller enters the next planning round with a slight upward bias on the breach-relevant capability dimensions. This coupling does not require retraining the adapter or modifying the case bank schema, and it means the pipeline's two auxiliary modules now share information through a common capability representation.

A noteworthy observation from the cross-scenario Reflection-only experiments reinforces the case for tight coupling: Reflection-only underperforms the Pure LLM baseline in all five generalization scenes (mean MS deficit $-0.0272$, single seed 7). This observation suggests that LLM-based self-critique, when applied in isolation without Memento's case-grounded evidence and Hope's structured weight adaptation, may not provide a reliable improvement signal and can actively misdirect the planner under the current setting. Multi-seed replication remains necessary to confirm whether this finding is systematic or seed-dependent. The implication is that reflection in this architecture functions as a diagnostic and repair signal that depends on the other components for grounding: Memento supplies evidence about what has worked in similar past cases, Hope ensures that the capability weights reflect both doctrinal priors and scenario-specific adaptation, and reflection then fine-tunes the plan against the simulator's feedback. Removing the grounding scaffolding leaves the critique untethered from actionable constraints.

The hyperparameter sensitivity analysis (Section~4.11) reveals a practically important finding: both HOPE SDE parameters ($\theta$ and $\epsilon$) exhibit U-shaped sensitivity curves with their reported default values near the performance minimum. The optimized Full pipeline ($\theta=0.9$, $\epsilon=0.005$, MS=0.6707) confirms that the HOPE component's potential is substantially underestimated at default parameters (original Full MS=0.6388 at $\theta=0.5$, $\epsilon=0.02$). The +0.0319 gain from hyperparameter optimization alone rivals or exceeds the gain from adding auxiliary components (memory, reflection) at default parameters. This finding connects to the case-bank size result: both ablations point toward the importance of systematic experimental design over post-hoc parameter selection.

The practical consequence is that the paper should be read not only as a claim of average improvement, but also as a map of where the present pipeline still appears structurally fragile. The multi-seed replication across five seeds demonstrates the Full pipeline's gain is consistent and statistically significant under default HOPE parameters (paired $t(4)=5.593$, $p<0.01$, Cohen's $d=2.50$), with the largest gains in seeds where the baseline is weakest. A supplementary multi-seed run with optimized parameters ($\theta=0.9$, $\epsilon=0.005$, $n=3$ seeds) confirms that the optimized Full outperforms Pure LLM in all three seeds (mean $\Delta$MS = +0.0606, $\text{SD}=0.0311$), extending the statistical significance finding to the best-performing configuration. However, seed 123 shows a reduced gain (+0.0262) compared to the default-parameter Full at the same seed (+0.0497), indicating that parameter optimization interacts with seed-specific initialization (Section~\ref{sec:results}, Table~\ref{tab:multiseed-optimized}). The 20-iteration convergence run confirms that extended optimization provides further improvement (+0.0869 MS over the 10-iteration result), but breach remains the persistent bottleneck. Just as importantly, the present loop should not be mistaken for a full reinforcement-learning or search regime in the style of standard RL formulations, AlphaGo-scale tree search, or population-based training \citep{sutton2018rl,silver2016alphago,jaderberg2017pbt}. It is closer to a bounded local update rule operating over multiple partially competing mission metrics, a setting that shares some flavor with multi-task optimization trade-offs \citep{yu2020pcgrad}.

\paragraph{Threats to validity.} The main quantitative gains are now more substantial after hyperparameter optimization. In the fixed-seed single-scenario ablation, the optimized Full-versus-Pure mission-success gain is +0.0447 (up from +0.0128 at default parameters). In the five-scene transfer suite, the scene-level mission-success gains range from +0.0104 to +0.0515. These differences are directionally encouraging, but readers should interpret the reported confidence intervals with care: all CIs in the cross-scenario tables reflect within-seed Monte Carlo variance only (50 runs, seed 7). The cross-seed standard deviation of Pure LLM mission success ($\pm0.0104$, Table~\ref{tab:multiseed}) is approximately 5$\times$ wider than the within-seed CI widths ($\sim$0.001--0.002), so CIs should not be interpreted as across-seed stability bounds. The multi-seed replication covers only Pure and Full at default parameters and should be extended to cover the optimized Full setting, per-component variants, and the Reflection-only cross-scenario result.

Several limitations remain before submission:

\begin{itemize}
\item \textbf{LLM backend traceability}: the refreshed canonical ablation and 20-iteration scene-out evidence now serialize machine-readable runtime manifests. The rerun records identify the local backend model as \path{deepseek-r1-0528-qwen3-8b} and place it in the Python virtual environment \texttt{Memento} on a 13th Gen Intel(R) Core(TM) i7-13700F CPU with an NVIDIA GeForce RTX 3060 GPU, using an OpenAI-compatible local endpoint. This removes the earlier dependence on author-confirmed runtime metadata for the main reported bundles.
\item \textbf{Experimental depth}: the multi-seed replication across five seeds ($\{7, 42, 123, 999, 2024\}$) now provides paired $t$-test evidence ($t(4)=5.593$, $p<0.01$) that the Full pipeline's improvement over Pure LLM is statistically significant and robust to seed variation. The three external baselines (few-shot CoT, random search, single-pass Memento) provide meaningful lower and intermediate bounds. The remaining limitation is that the multi-seed evidence currently covers only Pure and Full settings; the per-component ablation (Memento-only, Hope-only, Reflection-only) is available only for seeds 7 and 42. Additionally, the evidence still comes from a single scenario family.
	\item \textbf{Reflection-only cross-scenario result}: the finding that Reflection-only underperforms Pure LLM in all five generalization scenes (mean deficit $-0.0272$) is currently observed under a single seed (seed 7). Multi-seed replication is needed to confirm whether this negative result is systematic or seed-dependent.
\item \textbf{Case-bank comparability}: the ablation evidence uses a five-record frozen seed bank, whereas the selected generalization evidence uses 250 source cases and 200-case scene-out banks. This is informative, but it also means that the role of memory quality and memory scale must be discussed carefully. In particular, the weak Memento-only result in the ablation table may reflect severe case-bank undersizing (five records) as much as memory-module weakness. A case-bank-size ablation (5/25/100/200 records) is reported in Section~4.10 and confirms that memory scale matters: the 5-record frozen seed bank is indeed undersized, and the 100-record setting achieves the best overall result.
\item \textbf{Runtime-cost reporting}: the refreshed result bundles now record runtime manifests and wall-clock timing summaries. The cost--quality table provides approximate wall-clock estimates for all settings and baselines. Precise timing instrumentation will be added for the camera-ready version.
\item \textbf{Figure scope}: the main-text figures are now generated directly from repository-local JSON files, and a separate graphical abstract has been prepared as a non-generative submission asset. The remaining work is narrower: final grayscale verification, any journal-specific artwork packaging, and any extra venue-driven figure refinements.
\item \textbf{Metadata completeness}: the author list, affiliation, corresponding-author email, no-funding statement, competing-interest declaration, and CRediT roles are now recorded, but the final acknowledgments wording and data-availability wording may still need journal-specific refinement.
\item \textbf{Build readiness}: the manuscript now compiles successfully with a workspace-local TinyTeX toolchain, but the final submission package still needs the remaining editorial metadata and any journal-specific packaging checks.
\end{itemize}

These limitations do not invalidate the draft; they define the next cleanup stage. The manuscript is already useful as a structured, source-traceable application paper skeleton. The multi-seed replication ($n=5$, $t(4)=5.593$, $p<0.01$) provides strong statistical evidence for the Full-pipeline gain, and the external baselines bracket its performance between an undirected lower bound (few-shot CoT at MS=0.5189) and the optimized Full upper bound (MS=0.6707). The 20-iteration convergence analysis and BottleneckDetector diagnostics provide a mechanistic account of where and why the pipeline improves, and where it remains fragile. The next iteration should focus on multi-seed replication of the optimized Full setting, per-component variants, and Reflection-only cross-scenario results, as well as extending the evidence to additional scenario families.
